\chapter{Wstęp}

Dynamiczny rozwój esportu sprawił, że profesjonalne rozgrywki w grach komputerowych stały się
pełnoprawnym obszarem analizy ilościowej. Jednym z najgłośniejszych przykładów jest \textit{League
of Legends}, dla którego dostępne są wieloletnie dane o meczach, składach, statystykach gry oraz --
dla części spotkań -- kursach bukmacherskich. Taka kombinacja danych pozwala badać nie tylko to, czy
model potrafi wskazać zwycięzcę, lecz także czy potrafi przypisać realistyczne prawdopodobieństwo
wyniku.

Prognozowanie wyników takich meczów jest trudne, ponieważ obserwowana siła zespołu zmienia się w
czasie. Składy są przebudowywane, zawodnicy pełnią wyspecjalizowane role, a aktualizacje gry
wpływają na dominujące strategie. Dlatego sama historyczna nazwa drużyny jest często zbyt ubogą
reprezentacją jej bieżącego potencjału sportowego. W pracy przyjęto więc perspektywę przedmeczowej
predykcji probabilistycznej, w której model może korzystać wyłącznie z informacji dostępnych przed
rozpoczęciem spotkania.

Celem niniejszej pracy jest opracowanie oraz ocena modelu przewidującego wynik profesjonalnego meczu
\textit{League of Legends}. Model łączy systemy rankingowe zawodników, miary niepewności, pełny
historyczny kontekst drużynowy W20 oraz korektę dwumianową zależną od formatu serii Bo1/Bo3/Bo5.
Jako zewnętrzny punkt odniesienia wykorzystano prawdopodobieństwa implikowane przez kursy
bukmacherskie, które są jednym z najsilniejszych publicznie dostępnych benchmarków dla predykcji
przedmeczowej.

Praca nie dotyczy budowy systemu inwestycyjnego ani oceny opłacalności zakładów. Analiza ogranicza
się do jakości predykcji probabilistycznej, dlatego wykorzystano AUC, LogLoss, Brier Score oraz
Expected Calibration Error, przy czym główną metryką jest LogLoss.

Zakres pracy obejmuje analizę danych GOL.GG, przygotowanie wspólnej próby z kursami OddsPortal,
porównanie systemów rankingowych, budowę metamodelu W20-Binomial oraz końcową ocenę względem rynku
bukmacherskiego. Główna ewaluacja została przeprowadzona na meczach, dla których dostępne były oba
typy danych.

Główne pytanie badawcze pracy brzmi: czy model oparty na rankingach zawodników i historycznym
kontekście drużyn potrafi skutecznie przewidywać wyniki profesjonalnych meczów \textit{League of
Legends} oraz osiągać jakość probabilistyczną konkurencyjną względem kursów bukmacherskich?

W pracy przyjęto następujące hipotezy badawcze:

\begin{enumerate}
    \item Systemy rankingowe oparte na zawodnikach osiągają lepszą jakość predykcji niż systemy rankingowe oparte wyłącznie na drużynach.
    \item Metamodel łączący sygnały rankingowe poprawia jakość predykcji względem pojedynczego najlepszego systemu rankingowego.
    \item Dodanie historycznego kontekstu formy drużyn poprawia jakość probabilistyczną modelu względem wariantu opartego wyłącznie na rankingach.
    \item Model rankingowo-kontekstowy osiąga wyniki konkurencyjne względem prawdopodobieństw implikowanych przez kursy bukmacherskie.
\end{enumerate}

Dalsza część pracy została zorganizowana w następujący sposób. Najpierw przedstawiono stan wiedzy i
kontekst problemu, obejmujący predykcję sportową, systemy rankingowe, metamodelowanie, ocenę
probabilistyczną oraz rolę kursów bukmacherskich jako benchmarku. Następnie opisano środowisko
\textit{League of Legends}, wykorzystane źródła danych GOL.GG i OddsPortal oraz proces konstrukcji
wspólnej próby badawczej. Kolejny rozdział zawiera eksploracyjną analizę danych, obejmującą
strukturę zbioru, formaty meczów, poziomy rozgrywek oraz podstawową analizę rynku bukmacherskiego. W
dalszej części opisano metodykę ewaluacji, systemy rankingowe wykorzystane jako modele bazowe,
konstrukcję metamodelu W20-Binomial, dobór okna historycznego, porównanie rodzin estymatorów oraz
końcową symetryzację i kalibrację. Rozdział wynikowy prezentuje końcowe porównanie modeli, analizę
bootstrapową, segmentację według formatu serii i weryfikację hipotez. Następnie omówiono
interpretację wyników, ograniczenia badania oraz kierunki dalszych prac. Ostatni rozdział zawiera
syntetyczne podsumowanie i odpowiedź na pytanie badawcze.
